\section{Source and delivery note}
For an adult parish assembly: Twenty-ninth Sunday in Ordinary Time, Year A, 18 October 2026. Roman Missal, Third Edition, U.S. English (2011), Week XXIX; U.S. Lectionary for Mass, second typical edition, volume I, no.145. No local-calendar overlay or Communion-antiphon selection is assumed. Mark is named as one Communion alternative, with its surrounding Gospel context distinguished.

The argument joins the reviewed studies' ordered allegiance to grace becoming willing service: Aquinas's divine-image reading reaches the person without erasing Chrysostom's genuine civil dues or their moral limit. Cyrus and Paul's thanksgiving develop the return of that person through grateful service. The weekly response is an original pastoral application.

Spoken body: 1,302 words; approximately 10.0--11.3 minutes at 115--130 words per minute. This estimate is not a timed human delivery. Silent rereading checked sense, transitions, sentence length and oral clarity; no audible or human rehearsal is claimed. The speech ends before this note.

Scripture is explained in original paraphrase from the public-domain Douay--Rheims Bible, Challoner revision; the U.S. Lectionary controls proclamation. Missal prayers are explained, not translated for recitation. This is a study resource, not an official text or ecclesiastical approval. The canonical study retains the evidence limits, including uninspected exact 2008 Latin and publisher-specific 2011 altar-book pages.

\section{References}
\begin{description}[style=nextline,leftmargin=0pt,labelsep=0pt,itemsep=.35em]
\item[Scripture and prayer]
Matthew 22:15--21; Isaiah 45:1,4--6, with 44:26--28 as restoration context; Psalm 96:1,3--5,7--10a,c (Hebrew; Vulgate 95); 1 Thessalonians 1:1--5b; Mark 10:32--45 as context for the optional Communion excerpt at 10:45. Douay--Rheims/Challoner study text, Project Gutenberg 1581. Roman Missal, Week XXIX: Collect, Prayer over the Offerings, Mark Communion alternative and Prayer after Communion; editions and use identified above.
\item[John Chrysostom]
\work{Homilies on Matthew}, LXX, §§1--2, tribute and godliness; LXV, §§4--5, Christ's self-gift against ambition, on Matthew 20:28, parallel to Mark. Prevost translation revised by M.B. Riddle, NPNF I.10 (1888). \work{Homilies on 1 Thessalonians}, I, on 1:1--7, thanksgiving and active virtues; Oxford translation revised for NPNF I.13 (1889), CCEL. Verse 6 is unappointed context.
\item[Thomas Aquinas and Jerome]
Aquinas, \work{Super Matthaeum}, 22, lectio 2, Corpus Thomisticum paragraph 87389, tribute and the divine image; \work{Super Philippenses}, 2, lectio 3, paragraph 87838, divine action within willing. Jerome, \work{Commentarii in Isaiam}, XII on 45:1--7, DCO Latin \work{Pars 3}, pp.126--129, Cyrus and the gifted teacher. These are electronic witnesses; DCO's historical exemplar remains unestablished. Arguments are paraphrased, not quoted.
\end{description}
